\chapter{相关背景}

\section{检索增强生成}

\subsection{基本原理}

检索增强生成（Retrieval-Augmented Generation，RAG）是一种将外部知识检索
与大语言模型生成相结合的技术框架。
其核心思想是：在推理阶段，先从外部知识库中检索与用户查询最相关的文档片段，
再将这些片段作为上下文注入提示（Prompt），由语言模型基于上下文生成最终回复。

与直接依赖模型参数内隐知识的生成方式相比，RAG具有以下优势：
第一，模型无需记忆全部领域知识，知识库可独立更新而无需重新训练模型；
第二，通过控制检索到的文档质量，可有效降低模型幻觉的发生概率；
第三，检索来源可追溯，有助于提升系统输出的可解释性。

\subsection{RAG系统的基本组成}

一个典型的RAG系统由以下三个模块组成：

\begin{enumerate}
    \item \textbf{知识库构建模块}：将原始文档进行分块（Chunking）、
    编码（Encoding）并存储至向量数据库，形成可检索的索引。
    分块策略（大小、重叠长度）对检索质量有重要影响。

    \item \textbf{检索模块}：接收用户查询，将其编码为向量，
    在向量数据库中查找语义相近的文档块，并按相关性排序返回Top-$k$结果。

    \item \textbf{生成模块}：将检索到的文档块拼接至提示模板，
    调用大语言模型生成最终回复。
\end{enumerate}

\subsection{RAG技术的发展脉络}

RAG技术的研究起源于对大语言模型知识局限性的系统性认识。
早期的语言模型依赖固定参数内隐知识，一旦训练完成便无法获取新知识，
在时效性强或高度专业化的问答场景中表现明显不足。
Lewis等人在2020年提出的原始RAG框架将文档检索器与序列生成模型端到端联合训练，
首次验证了"检索+生成"范式在开放域问答任务上的有效性。

此后，RAG研究沿两条主线演进：
其一是\textbf{检索精度提升}，通过双编码器微调、交叉编码器重排序\citing{nogueira2019reranking}、迭代检索等手段，
使检索模块返回的上下文文档与查询意图更加吻合；
其二是\textbf{系统架构优化}，从最初的一次性检索扩展至多步推理中的多轮检索，
并引入检索必要性判断（即"是否应检索"的路由决策），降低无关文档引入的噪声。

面向专业领域的RAG系统还需额外应对\textbf{词汇鸿沟}与\textbf{知识库覆盖不足}两个问题\citing{gao2023ragsurvey}。
词汇鸿沟指用户口语化查询与专业文档之间的表述差异，
覆盖不足则指领域知识库规模有限时检索召回率低的问题。
面向领域的RAG工程实践通常通过同义问法注入、失败查询覆盖等启发式增强手段
缓解词汇鸿沟问题，并以评估驱动的迭代流程持续改善知识库质量，
这已成为当前领域RAG系统构建的主流范式之一。

\section{向量检索与嵌入模型}

\subsection{稠密向量检索}

稠密向量检索（Dense Retrieval）是RAG系统的核心检索机制。
其基本流程为：利用嵌入模型（Embedding Model）将文本映射至高维稠密向量空间，
查询文本与文档文本在同一向量空间中通过余弦相似度（Cosine Similarity）衡量语义相关性。
对于查询向量$\boldsymbol{q}$与文档向量$\boldsymbol{d}$，余弦相似度定义为：

\begin{equation}
\text{sim}(\boldsymbol{q}, \boldsymbol{d}) = \frac{\boldsymbol{q} \cdot \boldsymbol{d}}{\|\boldsymbol{q}\| \cdot \|\boldsymbol{d}\|}
\end{equation}

经归一化后，向量数据库中存储的余弦距离$dist$与相似度分数之间满足：

\begin{equation}
\text{score}_\text{vec} = \frac{1}{1 + dist}
\end{equation}

面向中文场景的嵌入模型需在中文语义相似度任务上表现优异，
同时参数量适中以支持本地部署。
BAAI发布的bge系列模型\citing{xiao2023cpack}针对中文文本优化，
在多个中文语义基准上取得较好结果，是中文RAG系统的常见选择之一。

\subsection{ChromaDB向量数据库}

向量数据库是RAG系统持久化存储文档向量的核心组件，
通常支持向量相似度查询、元数据过滤与集合管理等功能。
ChromaDB是一款开源的本地向量数据库，支持持久化存储与元数据过滤，
可在无网络依赖的环境中运行，适合隐私敏感或离线部署场景。
元数据过滤能力使系统在检索时能够将候选范围限制在特定类别文档中，
从而提升检索精度。

\subsection{文档分块策略}

一种常见的领域文档分块策略是Markdown感知的自适应分块：
优先按Markdown标题和双换行符切分段落以保留语义边界，
在段落粒度内控制块大小和重叠长度，
对超长段落回退至字符级滑动窗口切分。
此策略能够保留文档的语义结构，避免将同一语义单元割裂到不同块中。

\section{混合检索与重排序}

\subsection{稀疏检索的互补价值}

纯向量检索在语义相似度匹配上表现良好，
但对精确术语（如命令名称\texttt{tar}、\texttt{chmod}等）的检索存在不足：
若查询中出现罕见词汇而嵌入模型对其语义泛化处理，
可能导致语义正确但关键词不匹配的文档排名靠后。
稀疏检索方法（以BM25为代表\citing{robertson2009bm25}）基于词频统计，
在精确关键词匹配上具有天然优势，与向量检索形成互补。
实际工程中也存在更轻量的替代方案：
以查询词覆盖率（即命中词数与查询词数之比）作为词法相关性评分，
无需维护全局词频统计，实现更简单，
在小规模领域知识库上仍能提供有效的关键词补强信号。

\subsection{混合检索的实现方式}

混合检索的一种常见实现方式是以加权乘积将向量相似度与关键词命中率合并，
引入关键词命中率作为辅助信号对候选文档进行重排。
设文档$d$的向量相似度分数为$s_\text{vec}(d)$，关键词命中率为$s_\text{kw}(d)$，
则混合重排分数可定义为：

\begin{equation}
\text{score}(d) = s_\text{vec}(d) \times \bigl(1.0 + \alpha \cdot s_\text{kw}(d)\bigr)
\end{equation}

其中$\alpha$为关键词权重系数，通过在评估集上调参确定。
乘积形式保证了向量分数低的文档即使关键词完全命中也不会被过度提升，
从而在语义相关性与关键词精确性之间取得平衡。
此外，为防止低质量关键词候选引入噪声，通常设置命中率阈值对候选集进行过滤。

\subsection{重排序流程}

混合检索重排序的完整流程如下：
\begin{enumerate}
    \item 分别构建向量候选集与关键词候选集；
    \item 以文档块ID为键合并两个候选集，
    向量候选已有记录时直接补充其关键词分数，
    纯关键词候选通过命中率阈值过滤后加入合并集合；
    \item 按混合分数降序排列，截取Top-$k$返回。
\end{enumerate}

\section{Shell命令安全}

\subsection{Shell命令执行的风险特征}

Shell命令执行具有以下风险特征，使其成为LLM辅助系统的高风险组件：

\begin{enumerate}
    \item \textbf{强副作用性}：文件删除（\texttt{rm}）、权限变更（\texttt{chmod}）、
    磁盘写操作（\texttt{dd}）等命令一旦执行，影响往往不可逆；

    \item \textbf{组合爆炸性}：Shell命令通过管道、重定向等操作符组合，
    单个操作符的语义改变可能导致截然不同的执行结果；

    \item \textbf{提示注入风险}：若用户输入中包含精心构造的自然语言，
    可能诱导LLM生成与意图不符甚至危险的命令
    （即提示注入攻击，Prompt Injection）。
\end{enumerate}

\subsection{命令安全的防护层次}

针对上述风险，业界通常从以下几个层次构建Shell命令安全防护：

\textbf{输出过滤层}：在LLM生成命令后，通过规则引擎或分类模型检测危险模式，
拦截高风险命令。

\textbf{权限控制层}：基于白名单/黑名单机制，
明确定义允许执行的命令范围，对超出范围的命令予以拒绝。

\textbf{预执行模拟层}：在实际执行前以只读或模拟方式预览命令效果，
使用户能够在确认结果后再决定是否执行。

本系统的安全框架覆盖以上全部三个层次，
具体实现详见第五章。

\section{模型量化技术}

\subsection{大模型量化的必要性}

大语言模型的参数量通常以十亿（Billion）为单位，
以FP32精度存储时，7B参数模型约需28 GB内存，
超出消费级GPU（通常4--16 GB VRAM）的承载能力。
量化（Quantization）技术通过降低参数的数值精度
将模型压缩至可在消费级硬件上运行的规模，
是大模型本地化部署的核心使能技术。

量化方法可按压缩时机分为两类：
\textbf{训练后量化（Post-Training Quantization，PTQ）}在训练完成后直接对权重进行精度转换，
无需重新训练，适合快速部署；
\textbf{量化感知训练（Quantization-Aware Training，QAT）}在训练阶段引入量化误差反馈，
精度损失更小但需要额外的计算资源。
本系统采用PTQ方案，基于Hugging Face Transformers集成的BitsAndBytes库实现。

\subsection{NF4量化原理}

BitsAndBytes库提供的4-bit量化方案支持两种数值格式：
\textbf{FP4}（Float 4-bit）和\textbf{NF4}（Normal Float 4-bit）。
NF4由Dettmers等人在QLoRA工作中提出\citing{dettmers2023qlora}，
其核心思想是：LLM的权重在经过适当归一化后近似服从零均值正态分布，
因此可以设计一种信息论意义上近优的4-bit数字表示，
使量化级别在正态分布高概率密度区域分布更密集，
从而降低量化误差的期望值。
相比FP4的均匀量化，NF4在同等比特宽度下能够以更少的量化误差表示正态分布的权重，
是当前4-bit量化的主流选择。

\subsection{双重量化与本系统配置}

为进一步压缩内存开销，本系统同时启用\textbf{双重量化（Double Quantization）}：
对存储量化常数（scaling factors）所用的FP32精度本身再次进行8-bit量化，
平均每个参数额外节省约0.37 bit的存储空间。

Hugging Face BitsAndBytes库将NF4与双重量化集成为开箱即用的接口，
支持在加载模型时通过统一配置同时启用4-bit量化、双重量化与计算精度设置，
矩阵乘法阶段通常以BF16或FP16精度执行以兼顾速度与数值稳定性。
Transformers的\texttt{device\_map}机制可将模型权重在GPU与CPU内存之间自动均衡分配，
使超出单卡VRAM容量的模型得以在消费级硬件上完成推理。
本文采用上述技术组合部署大语言模型，具体参数配置详见第三章系统设计。

\section{本章小结}

本章介绍了与本系统相关的核心技术背景。
首先阐述了检索增强生成（RAG）的基本原理、系统组成与技术发展脉络；
其次介绍了稠密向量检索的数学基础、中文嵌入模型的选型背景及文档分块策略；
然后分析了稀疏检索与向量检索的互补性，并介绍了基于加权乘积的混合检索重排方法；
再从风险特征与防护层次两个角度概述了Shell命令安全的技术背景；
最后介绍了大模型量化技术的基本原理与NF4、双重量化方案的核心思路。
这些技术构成了本系统设计与实现的理论基础，
后续章节将在此基础上展开具体的系统设计与实验分析。
